% ============================================================
% Supplementary Material for IROS 2026 Submission
% ============================================================
\documentclass[letterpaper, 10pt]{article}
\usepackage[utf8]{inputenc}
\usepackage{amsmath,amssymb}
\usepackage{graphicx}
\usepackage{booktabs}
\usepackage{hyperref}
\usepackage{geometry}
\usepackage{enumitem}
\usepackage{xcolor}
\usepackage{subcaption}
\usepackage{pifont}
\usepackage{times}

\geometry{margin=1in}

\newcommand{\sysname}{Agentic RAG-VLM}
\newcommand{\haarag}{HAA-RAG}
\newcommand{\cmark}{\ding{51}}
\newcommand{\xmark}{\ding{55}}

\title{Supplementary Material:\\Agentic RAG-VLM: Affordance-Aware Retrieval-Augmented Generation\\with Self-Reflective Planning for Robotic Grasping}

\author{Anonymous IROS Submission}

\begin{document}
\maketitle

\section{Extended Ablation Results}

\subsection{Agentic Pipeline Component Analysis}

Table~\ref{tab:agentic_ablation} shows the contribution of each agentic component, added incrementally to the VLM-Only baseline.

\begin{table}[h]
\centering
\caption{Agentic pipeline ablation. Each row adds one component cumulatively.}
\label{tab:agentic_ablation}
\begin{tabular}{@{}lccc@{}}
\toprule
\textbf{Configuration} & \textbf{Success (\%)} & \textbf{Avg Retries} & $\Delta$ \\
\midrule
A. VLM-Only & 25.0 & 0.00 & -- \\
B. + HAA-RAG & 45.8 & 0.00 & +20.8 \\
C. + Retry (3-level) & 67.2 & 1.12 & +42.2 \\
D. + Scene Graph & 68.1 & 1.08 & +43.1 \\
E. + Episodic Memory & \textbf{69.7} & 1.04 & \textbf{+44.7} \\
\bottomrule
\end{tabular}
\end{table}

\textbf{Analysis.} RAG retrieval provides the largest single improvement ($+20.8\%$), enabling the system to leverage stored affordance-specific grasp parameters rather than generic VLM heuristics.
The three-level retry mechanism adds another $+21.4\%$, confirming that closed-loop failure recovery is critical for interactive and long-horizon tasks.
Scene graph constraints and episodic memory each contribute incremental gains of $+0.9\%$ and $+1.6\%$, respectively.
The memory contribution is concentrated on long-horizon tasks (51.7\% vs.\ 36.7\% without memory), where parameter transfer between similar objects is most effective.

\subsection{Per-Task Success Breakdown}

\begin{table}[h]
\centering
\caption{Per-task success rates (\%) for the full system.}
\label{tab:per_task}
\setlength{\tabcolsep}{3pt}
\begin{tabular}{@{}llcc@{}}
\toprule
\textbf{Category} & \textbf{Task} & \textbf{Success} & \textbf{Avg Retries} \\
\midrule
\multirow{6}{*}{Single} & Pick cup & 100.0 & 0.0 \\
 & Pick apple & 100.0 & 0.0 \\
 & Pick bottle & 96.7 & 0.3 \\
 & Pick banana & 100.0 & 0.0 \\
 & Pick bowl (edge grasp) & 56.7 & 1.8 \\
 & Pick cube & 50.0 & 2.1 \\
\midrule
\multirow{4}{*}{Interactive} & Cup near vase & 60.0 & 1.5 \\
 & Ball behind box & 76.7 & 0.8 \\
 & Phone near glass & 53.3 & 1.9 \\
 & Pick with tool & 40.0 & 2.3 \\
\midrule
\multirow{2}{*}{Long-horizon} & Clear table (3 objects) & 53.3 & 1.2 \\
 & Sort by fragility & 50.0 & 1.6 \\
\bottomrule
\end{tabular}
\end{table}

\textbf{Analysis.}
Easy single-grasp tasks (cup, apple, banana) achieve 100\% success with zero retries, demonstrating reliable RAG retrieval for well-represented affordance types.
Bowl and cube tasks are more challenging due to edge-grasp and geometric constraints.
Interactive tasks range from 40--77\% depending on constraint complexity.
Long-horizon tasks hover around 50\%, limited by error accumulation across sequential grasps.

\section{Failure Type Taxonomy}

Table~\ref{tab:failure_types} lists the 14 failure types organized in four groups, with associated detection signals and recovery strategies.

\begin{table}[h]
\centering
\caption{Failure type taxonomy and recovery strategies.}
\label{tab:failure_types}
\setlength{\tabcolsep}{3pt}
\begin{tabular}{@{}llll@{}}
\toprule
\textbf{Group} & \textbf{Failure Type} & \textbf{Detection} & \textbf{Recovery} \\
\midrule
\multirow{3}{*}{Pre-contact}
 & POSITION\_ERROR & IK residual $> \epsilon$ & Re-localize target \\
 & UNREACHABLE & IK divergence & Adjust approach dir. \\
 & APPROACH\_ERROR & Angle mismatch & Rotate approach $\pm$15$^\circ$ \\
\midrule
\multirow{3}{*}{Contact}
 & WIDTH\_MISMATCH & Gripper fails close & $w' = w \pm 10$\,mm \\
 & COLLISION & Non-target contact & Height $+$30\,mm \\
 & ORIENTATION\_ERROR & Grasp angle $>$ tol. & Rotate 90$^\circ$ \\
\midrule
\multirow{4}{*}{Post-grasp}
 & SLIP & Force $\downarrow$ during lift & $f' = 1.2f$ \\
 & DROP & Object falls post-lift & $f' = 1.3f$, verify \\
 & FORCE\_DAMAGE & $f > f_{\max}^{\text{frag}}$ & $f' = 0.7f$ \\
 & DEFORMATION & $\Delta$ shape detected & Switch to pinch \\
\midrule
\multirow{4}{*}{Task-level}
 & WRONG\_OBJECT & ID mismatch & Re-identify target \\
 & CONSTRAINT\_VIOL. & SG constraint fail & Apply constraints \\
 & TIMEOUT & $t > t_{\max}$ & Simplify plan \\
 & UNKNOWN & Unclassified & Full re-planning \\
\bottomrule
\end{tabular}
\end{table}

\section{Knowledge Base Statistics}

Our knowledge base contains 110 curated grasp experiences spanning 10 object categories and 8 affordance types.

\begin{table}[h]
\centering
\caption{Knowledge base composition (110 experiences, 8 affordance types).}
\begin{tabular}{@{}lccl@{}}
\toprule
\textbf{Category} & \textbf{Count} & \textbf{Success \%} & \textbf{Primary Affordance} \\
\midrule
Cup/Mug & 15 & 87\% & GRASPABLE\_HANDLE \\
Bottle & 12 & 83\% & GRASPABLE\_BODY \\
Bowl & 10 & 70\% & GRASPABLE\_EDGE \\
Fruit (apple, banana) & 14 & 93\% & WRAPPABLE \\
Cube/Block & 10 & 80\% & GRASPABLE\_BODY \\
Glass (vase, wine) & 12 & 67\% & FRAGILE \\
Tool (screwdriver) & 10 & 75\% & TOOL\_HANDLE \\
Phone/Electronics & 10 & 70\% & CLAMPABLE \\
Ball & 8 & 88\% & WRAPPABLE \\
Other & 9 & 78\% & Various \\
\midrule
\textbf{Total} & \textbf{110} & \textbf{79\%} & \textbf{8 types} \\
\bottomrule
\end{tabular}
\end{table}

Each experience stores: object category, affordance descriptor (type, material, fragility, region), grasp parameters (position offset, width, force, grasp type), and binary success label.
Failed experiences are retained with contrastive penalties to prevent strategy repetition.

\section{Computational Cost Breakdown}

\begin{table}[h]
\centering
\caption{End-to-end latency breakdown per grasp decision (INT4 quantized).}
\begin{tabular}{@{}lccc@{}}
\toprule
\textbf{Stage} & \textbf{Latency} & \textbf{Share} & \textbf{Notes} \\
\midrule
Scene Analysis (VLM) & 850 ms & 45\% & INT4 quantized \\
RAG Retrieval (3-level) & 15 ms & 0.8\% & In-memory indexed \\
Scene Graph Construction & 50 ms & 2.7\% & Object-pair traversal \\
Constraint Inference & 10 ms & 0.5\% & 4 constraint types \\
Grasp Planning (VLM) & 800 ms & 42\% & Parameter generation \\
Reflection (if retry) & 150 ms & 8\% & Conditional, per retry \\
\midrule
\textbf{Total (first attempt)} & \textbf{$\sim$1.7 s} & -- & -- \\
\textbf{Total (with 1 retry)} & \textbf{$\sim$2.7 s} & -- & -- \\
\bottomrule
\end{tabular}
\end{table}

\section{Reproducibility}

\textbf{Code}: Available at \url{[anonymous upon acceptance]}.

\textbf{Hardware}: VLM inference benchmarked on NVIDIA RTX 5090 (32GB). The INT4-quantized model (6.7\,GB) can also run on NVIDIA Jetson AGX Orin.
Simulation experiments use CPU-based physics-informed simulation (no GPU required).

\textbf{Software}: Python 3.11, NumPy 1.26.4, SciPy 1.11.4.
VLM: Qwen3-VL-8B with BitsAndBytes INT4 quantization.

\textbf{Seeds}: All experiments use seed=42 with 30 trials per task.
Statistical testing uses 5 independent complete repetitions with two-tailed $t$-tests ($\alpha = 0.05$, Bonferroni-corrected).

\textbf{Simulation}: Physics-informed tabletop simulation modeling Franka Panda 7-DoF kinematics, parallel-jaw gripper (0--80\,mm), workspace geometry, IK feasibility, geometric collision detection, and seven-factor grasp quality evaluation.
Grasp outcomes are determined by structured threshold analysis on quality factors, ensuring reproducible and physically meaningful results.

\textbf{Video}: The supplementary video demonstrates single-object grasping, interactive grasping with collision avoidance, and failure recovery through self-reflection across representative tasks.

\end{document}
